\begin{abstract}
We introduce \textsc{QLoRA}, a novel finetuning method designed to drastically lower memory requirements, making it possible to finetune models with up to 65B parameters on a single 48GB GPU, all while retaining the accuracy of standard 16-bit finetuning for downstream tasks. The \textsc{QLoRA} technique enables gradient updates to be computed through a fixed, 4-bit quantized base language model into Low Rank Adapters~(LoRA). The resulting best-performing model suite, which we call \textbf{Guanaco}, surpasses all previously available open models on the Vicuna benchmark, achieving 99.3\% of ChatGPT’s benchmark performance and requiring only 24 hours of finetuning on a single GPU. \textsc{QLoRA} incorporates several key techniques to minimize memory usage while maintaining accuracy: (a) 4-bit NormalFloat (NF4), an innovative data format that achieves information-theoretic optimality for weights following a normal distribution; (b) Double Quantization, a method that further compresses memory use by quantizing the quantization parameters; and (c) Paged Optimizers, which mitigate memory spikes during training. Utilizing \textsc{QLoRA}, we finetuned over 1,000 models and conducted a comprehensive evaluation of instruction-following and chatbot quality on 8 instruction-tuning datasets, with different architectures (LLaMA, T5) and at scales that standard finetuning cannot accommodate (e.g., models with 33B and 65B parameters). Our findings indicate that applying QLoRA finetuning to a small, high-quality dataset delivers state-of-the-art results—even when leveraging models smaller than prior leading methods. We further provide extensive analyses of chatbot performance, using both human judgments and GPT-4-based evaluation, and demonstrate that the latter offers an inexpensive and fairly reliable substitute for human ratings. Additionally, we discover that existing chatbot benchmarks fail to reliably distinguish the true capabilities of these models. Through a focused lemon-picked analysis, we illustrate specific cases where \textbf{Guanaco} falls short relative to ChatGPT. All models and code are publicly released, including CUDA kernels supporting 4-bit training.\footnote{\url{https://github.com/artidoro/qlora} and \url{https://github.com/TimDettmers/bitsandbytes}}
\end{abstract}

\section{Introduction}

Adapting large language models (LLMs) through finetuning is a proven strategy for enhancing their capabilities \citep{min2021metaicl, wei2021finetuned, ouyang2022training, wang2022super, wang2022self, liu2022few} as well as instilling or eliminating particular behaviors \citep{ouyang2022training,askell2021general,bai2022training}. Yet, the high computational and memory costs associated with finetuning massive models render it impractical; for instance, standard 16-bit finetuning of a LLaMA model with 65B parameters~\citep{touvron2023llama} exceeds 780 GB of GPU memory demand. Although advanced quantization approaches can lower LLM memory usage during inference \citep{dettmers2022llm, dettmers2022case, frantar2022gptq, xiao2022smoothquant}, these strategies fail to provide the same benefits during training \citep{wortsman2023stable}. 

In this work, we show, for the first time, that finetuning a quantized 4-bit model can be accomplished with no loss in performance. The proposed method, \textsc{QLoRA}, introduces an innovative high-precision quantization scheme for pretrained models down to 4 bits, supplemented with a compact set of learnable Low-rank Adapter weights \citep{hu2021lora}, 

which are updated via gradient backpropagation through the quantized parameter space. 

\textsc{QLoRA} dramatically lowers the average GPU memory necessary to finetune a 65B parameter model from over 780GB to under 48GB, without increasing computational runtime or compromising accuracy relative to a fully finetuned 16-bit counterpart. This advance greatly democratizes access to LLM finetuning by enabling state-of-the-art models to be finetuned on a single GPU. Leveraging \textsc{QLoRA}, we develop the \textbf{Guanaco} model family. The second-best model in this series attains 97.8\% of ChatGPT’s performance on the Vicuna~\citep{vicuna2023} benchmark, training in less than 12 hours on a consumer-grade GPU. With a single high-end GPU in 24 hours, our largest version achieves 99.3\%, effectively matching ChatGPT on Vicuna. Upon deployment, our most compact \textbf{Guanaco} model (7B parameters) requires merely 5 GB of memory, yet surpasses a 26 GB Alpaca model by over 20 points on the Vicuna evaluation (Table~\ref{tbl:vicuna_eval}).

\textsc{QLoRA} presents several advancements aimed at minimizing memory requirements while maintaining model performance: (1) \textbf{4-bit NormalFloat}, a quantization datatype that is theoretically optimal for normally distributed values and demonstrates superior empirical performance compared to both 4-bit Integers and 4-bit Floats. (2) \textbf{Double Quantization}, a technique wherein the quantization constants themselves are quantized, leading to an average reduction of approximately 0.37 bits per parameter, which translates to a memory saving of around 3 GB in a 65B parameter model. (3) \textbf{Paged Optimizers}, which utilize NVIDIA's unified memory architecture to prevent transient memory surges during gradient checkpointing, especially for mini-batches with long sequence lengths. Integrating these advances results in an enhanced LoRA configuration, deploying adapters throughout all network layers and thereby nearly eliminating the accuracy compromises observed in previous approaches.

The efficiency afforded by \textsc{QLoRA} makes it possible to systematically investigate instruction finetuning and conversational agent performance at parameter scales infeasible with standard finetuning because of excessive memory usage. To this end, we train over 1,000 models, spanning multiple instruction tuning datasets, architectural variants, and model sizes ranging from 80M to 65B parameters. Our experiments demonstrate that \textsc{QLoRA} achieves parity with 16-bit training (\S\ref{sec:comp_to_std}) and is capable of training the state-of-the-art chatbot, \textbf{Guanaco} (\S\ref{sec:pushing}). In addition, we assess the behavior of the resulting models, finding, first, that data quality outweighs dataset scale—for example, the 9k-sample OASST1 dataset surpasses a much larger 450k-sample FLAN v2 (subsampled) dataset on chatbot evaluations, despite both being designed to enhance instruction generalization. Second, we show that superior performance on the Massive Multitask Language Understanding (MMLU) benchmark does not correlate with outstanding results on the Vicuna chatbot benchmark, and the converse is also true; this indicates that, for a given task, the appropriateness of a dataset has a greater impact than its size.

In addition, we conduct a comprehensive assessment of chatbot capabilities through dual evaluation by both human raters and GPT-4. Our methodology employs a tournament-style evaluation framework, wherein models face off in direct comparisons to generate optimal responses to provided prompts. The victor in each round is selected either by human judges or GPT-4. The accumulated tournament outcomes are synthesized into Elo scores~\citep{elo1967proposed, elo1978rating}, which serve as a metric for establishing the comparative ranking of chatbot models. Our findings indicate a high degree of concordance between GPT-4-based and human-driven assessments regarding model ranking; however, we also observe notable episodes of significant divergence. These results underscore the utility and limitations of model-based evaluation, which, despite offering a more resource-efficient alternative to human annotation, introduces certain ambiguities.

To complement the quantitative results of our chatbot benchmarks, we present a qualitative examination of the \textbf{Guanaco} models. This qualitative review elucidates specific examples of success and failure that are not evident from numerical evaluation alone.

We have made available all chatbot outputs accompanied by corresponding human and GPT-4 annotations to promote additional research. Our full codebase, including CUDA kernels, is open-sourced, and the proposed methods have been integrated with the Hugging Face transformers library \citep{wolf2019huggingface} for community use. We also distribute a set of adapters compatible with 7/13/33/65B model sizes, each trained on eight distinct instruction-following datasets, yielding a total of 32 open-source, fine-tuned models.

\section{Background}
\paragraph{Block-wise k-bit Quantization} Quantization involves mapping data from a representation of higher precision to one of lower precision, typically by reducing the number of bits per data element—such as converting from 32-bit floating-point values to 8-bit integer values. To fully utilize the limited range of the lower-precision data type, the input data is typically rescaled to fit within the bounds of the new type, often through normalization using the maximum absolute value in the dataset, which is frequently arranged as a tensor. For instance, when transforming a tensor of 32-bit floating-point numbers (FP32) into a signed 8-bit integer tensor (Int8) whose values span $[-127, 127]$:

\begin{equation}
    \mathbf{X}^{ \text{Int8}} = \text{round}\left( \frac{127}{{\text{absmax}}(\mathbf{X}^{{\text{FP32}}})}\mathbf{X}^{{\text{FP32}}} \right) = \text{round}(c^{\text{FP32}}\cdot \mathbf{X}^{{\text{FP32}}}),
\end{equation}

where $c$ denotes the {\em quantization constant} or {\em quantization scale}. The reverse operation, dequantization, reconstructs the original floating-point values as follows:

\begin{equation}
    \text{dequant}(c^{\text{FP32}}, \mathbf{X}^{\text{Int8}}) = \frac{\mathbf{X}^{\text{Int8}}}{c^{\text{FP32}}} = \mathbf{X}^{\text{FP32}}
\end{equation}

A key limitation of this method arises when the input tensor contains outliers, i.e., values of unusually high magnitude. In such cases, the quantization process becomes inefficient, as certain bit combinations—quantization bins—end up underutilized or entirely unused, with very few or no data points assigned to them. To address this, a widely adopted solution involves partitioning the input tensor into several blocks, each of which is quantized separately using its own quantization constant $c$. More formally, we partition the tensor $\mathbf{X}\in \mathbb{R}^{b\times h}$ into $n$ consecutive blocks, each of size $B$, by first flattening the tensor and then dividing the resultant sequence into ${n = ({b\times h} )/{B} }$ segments. Each segment is then independently quantized according to Equation~1, resulting in a quantized tensor along with $n$ distinct quantization constants $c_i$.

\paragraph{Low-rank Adapters} The Low-rank Adapter (LoRA) technique \citep{hu2021lora} offers a memory-efficient approach to finetuning by introducing a limited number of tunable parameters—referred to as adapters—while keeping the original model weights frozen. During stochastic gradient descent, gradients are computed and propagated through the static pretrained parameters, but only the adapter components are updated to minimize the loss. LoRA enhances a linear transformation by incorporating an additional low-rank, factorized projection. Specifically, given a mapping ${\mathbf{X} \mathbf{W} = \mathbf{Y}}$ with $\mathbf{X}\in  \mathbb{R}^{b\times h}$ and $\mathbf{W}\in  \mathbb{R}^{h\times o}$, LoRA modifies the output as follows:

\begin{equation}
    \mathbf{Y} = \mathbf{X}\mathbf{W} + s\mathbf{X}\mathbf{L}_1\mathbf{L}_2,
\end{equation}

where $\mathbf{L}_1\in  \mathbb{R}^{h\times r}$, $\mathbf{L}_2\in  \mathbb{R}^{r\times o}$, and $s$ denotes a scalar coefficient.

\paragraph{Memory Requirement of Parameter-Efficient Finetuning} The memory demands associated with LoRA during training, particularly regarding both the quantity and the scale of adapters deployed, warrant careful consideration. Given LoRA's notably small memory footprint, it is feasible to utilize a greater number of adapters to boost model performance while incurring only a marginal rise in overall memory consumption. Although LoRA falls under the category of Parameter Efficient Finetuning (PEFT) approaches, the bulk of memory consumption in LLM finetuning stems from activation gradients rather than from the parameters introduced by LoRA itself. For example, with a 7B LLaMA model trained on FLAN v2 using a batch size of 1 and employing LoRA parameters equal to the commonly adopted 0.2\% of the model's original weights\citep{hu2021lora,liu2022few}, the memory utilized by the LoRA input gradients amounts to 567 MB, whereas the LoRA parameters themselves occupy merely 26 MB. When gradient checkpointing~\citep{chen2016training} is applied, the memory needed for input gradients drops substantially to about 18 MB per sequence, yet this remains larger than the memory taken up by all LoRA weights combined. For context, a 4-bit base model requires 5,048 MB of memory. This comparison demonstrates not only the importance of gradient checkpointing in reducing memory consumption, but also indicates that further decreasing LoRA parameter size offers limited gains in memory savings. Consequently, adding more adapters is viable without substantially impacting the memory requirements during training (see Appendix~\ref{app:memory} for a comprehensive memory analysis). As will be elaborated, this flexibility is vital for achieving the performance level of full 16-bit precision models.

\section{\textsc{QLoRA} Finetuning}

\textsc{QLoRA} enables effective 4-bit finetuning through two novel techniques introduced in this work—namely, 4-bit NormalFloat (NF4) quantization and Double Quantization. Alongside these methods, we present Paged Optimizers, designed to address and mitigate the issue of out-of-memory errors during gradient checkpointing, which have historically posed challenges to finetuning large models on a single machine.

Within \textsc{QLoRA}, parameters are maintained using a low-precision storage type—typically 4-bit—while computation is carried out using BFloat16 precision. In practice, whenever a weight tensor under \textsc{QLoRA} is accessed, it undergoes dequantization to BFloat16 before being used for 16-bit matrix multiplications.

In the following, we elaborate on the specific elements of \textsc{QLoRA} and subsequently present its formal definition.

\paragraph{4-bit NormalFloat Quantization} The NormalFloat (NF) data type extends Quantile Quantization \citep{dettmers2022optimizers}, an information-theoretically optimal quantization scheme that guarantees each quantization bin is assigned an equal fraction of values from the source tensor. Quantile quantization estimates the quantiles of the input tensor by leveraging the empirical cumulative distribution function.

A significant drawback of quantile quantization, however, lies in the computational cost of quantile estimation. To overcome this, efficient quantile approximation strategies, such as SRAM quantiles~\citep{dettmers2022optimizers}, are typically used. These methods, though expedient, can introduce sizable quantization errors—especially for outliers, which often carry substantial importance.

If the distribution of the input tensors remains consistent except for a scaling constant, both the high computational overhead and approximation-induced quantization errors can be largely circumvented. Under such distributional stability, the quantiles for all tensors align, making exact quantile estimation much more tractable.

Pretrained neural network weights typically follow a zero-centered normal distribution characterized by a standard deviation $\sigma$ (refer to Appendix~\ref{app:normality}). By appropriately rescaling $\sigma$, all weights can be mapped to a single, fixed distribution that aligns with the numerical range of our chosen data type. In our case, this range is specified as $[-1, 1]$. Consequently, it is necessary to normalize both the quantiles corresponding to our data type and the neural network weights to fall within this interval.

To obtain the information-theoretically optimal data type for encoding zero-mean normal distributions with arbitrary $\sigma$ values within $[-1, 1]$, the following procedure is employed: (1) compute the $2^k+1$ quantiles of a standard normal distribution $N(0, 1)$ to construct a $k$-bit quantile quantization data type suitable for such distributions, (2) normalize these quantile values into the interval $[-1, 1]$, and (3) transform an input weight tensor by rescaling it to fit inside $[-1, 1]$ through absolute-max normalization.

After harmonizing the scales of the weight values and the data type, standard quantization can be applied. The third step effectively rescales the standard deviation of the weights to match that of the $k$-bit quantile data type. Explicitly, the $2^k$ representative values $q_i$ for the data type are determined by:

\begin{equation}
q_i  = \frac{1}{2}\left( Q_X\left(\frac{i}{2^k+1}\right) + Q_X\left(\frac{i+1}{2^k+1}\right)\right),
\end{equation}

where $Q_X(\cdot)$ denotes the quantile function corresponding to the standard normal distribution $N(0,1)$. A limitation of conventional symmetric $k$-bit quantization is the absence of an exact representation for zero, which is crucial for losslessly quantizing padding and other elements valued at zero. To address this and guarantee the inclusion of a discrete zeropoint at $0$ while utilizing all $2^k$ states in a $k$-bit data type, we propose an asymmetric construction. Specifically, we estimate quantiles $q_i$ separately over two intervals: allocating $2^{k-1}$ bins to the negative domain and $2^{k-1}+1$ bins to the positive domain. These quantile sets are then merged, and one of the duplicate zeros—present in both partitions—is removed. We refer to the resulting scheme as \emph{k-bit NormalFloat} (NFk), which assigns an equal expected count of elements to each quantization interval, thus rendering it information-theoretically optimal for zero-mean Gaussian-distributed data. For the precise definition of this data type, please consult Appendix~\ref{app:nf4}.

\paragraph{Double Quantization{}} 
We propose \emph{Double Quantization} (DQ) to further economize memory by quantizing the quantization constants themselves. While achieving accurate 4-bit quantization necessitates small block sizes~\citep{dettmers2022case}, the resultant storage for quantization constants can be substantial. For instance, using 32-bit quantization constants with a block size of 64 for $\mathbf{W}$ leads to an average overhead of $32/64 = 0.5$ bits per parameter. Double Quantization aims to reduce this overhead.

In particular, Double Quantization considers the initial set of quantization constants $c_2^{\text{FP32}}$ as inputs for a subsequent quantization step. This secondary process produces quantized constants $c_2^{\text{FP8}}$ along with an additional set of quantization constants $c_1^{\text{FP32}}$. For the second-stage quantization, we utilize 8-bit floating-point representations with a block size of 256, as empirical evidence from \citet{dettmers2022case} suggests that such 8-bit quantization incurs negligible performance loss. 
Given that the $c_2^{\text{FP32}}$ values are nonnegative, we subtract their mean before quantization to achieve a zero-centered distribution, thereby allowing symmetric quantization. On average, this method—using a block size of 64—lowers the memory required per parameter for storing quantization constants from $32/64 = 0.5$ bits down to $8/64 + 32/(64\cdot256) = 0.127$ bits, representing a per-parameter savings of 0.373 bits.

\paragraph{Paged Optimizers} leverage the NVIDIA unified memory~\footnote{\tiny \url{https://docs.nvidia.com/cuda/cuda-c-programming-guide}} capability, which facilitates seamless transfer of memory pages between CPU and GPU, ensuring correct GPU computation even in instances of GPU memory exhaustion. This mechanism functions analogously to conventional memory paging that occurs between system RAM and persistent storage. In our approach, we employ this feature to allocate paged memory for the optimizer state variables. As a result, when the GPU lacks sufficient memory, these states are transparently offloaded to CPU RAM and are subsequently reloaded into GPU memory whenever the optimizer's update phase requires them.

\paragraph{\textsc{QLoRA}.} Incorporating the elements outlined previously, we specify the formulation of \textsc{QLoRA} for an individual linear layer within a quantized base model equipped with a single LoRA adapter by:

\begin{equation}
    \mathbf{Y}^{\text{BF16}} =  \mathbf{X}^{\text{BF16}} \text{doubleDequant}(c_1^{\text{FP32}}, c_2^{\text{k-bit}}, \mathbf{W}^{\text{NF4}}) + \mathbf{X}^{\text{BF16}}\mathbf{L}^{\text{BF16}}_1\mathbf{L}^{\text{BF16}}_2,
\end{equation}

The doubleDequant$(\cdot)$ operation is defined as:

\begin{equation}
    \text{doubleDequant}(c_1^{\text{FP32}}, c_2^{\text{k-bit}}, \mathbf{W}^{\text{k-bit}}) = \text{dequant}(\text{dequant}(c_1^{\text{FP32}}, c_2^{\text{k-bit}}), \mathbf{W}^{\text{4bit}}) = \mathbf{W}^{\text{BF16}},
\end{equation}

For the quantized weights $\mathbf{W}$, we utilize NF4, and employ FP8 format for $c_2$.

To achieve finer quantization accuracy, a block size of 64 is chosen for $\mathbf{W}$, while $c_2$ uses a block size of 256 to optimize memory utilization.

During parameter optimization, it is only necessary to compute the gradients of the adapter weights, i.e., $\frac{\partial E}{\partial \mathbf{L}_i}$; gradients with respect to the 4-bit quantized weights, $\frac{\partial E}{\partial \mathbf{W}}$, are not required. Nonetheless, deriving $\frac{\partial E}{\partial \mathbf{L}_i}$ inherently involves the computation of $\frac{\partial \mathbf{X}}{\partial \mathbf{W}}$, which relies on equation~(5). This process requires dequantizing $\mathbf{W}^{\text{NF4}}$ from its storage representation to the computation type $\mathbf{W}^{\text{BF16}}$, thereby enabling the calculation of $\frac{\partial \mathbf{X}}{\partial \mathbf{W}}$ at BFloat16 numerical precision.

In essence, \textsc{QLoRA} utilizes a single data type for storage—typically 4-bit NormalFloat—and a separate data type for computation, specifically 16-bit BrainFloat. During training, the stored values are dequantized to the computation format for both the forward and backward passes; however, weight gradients are only calculated for the LoRA-specific parameters, which are maintained in the 16-bit BrainFloat format.

\section{QLoRA vs. Standard Finetuning}
\label{sec:comp_to_std}

Having outlined the operational principles of QLoRA and its notable memory efficiency for model finetuning, we now turn to the question of how its effectiveness compares to full-model finetuning approaches. We are also interested in dissecting the contributions of different QLoRA elements, such as the influence of NormalFloat4 in contrast to traditional Float4. The sections that follow detail our experimental investigations into these aspects.

\paragraph{Experimental setup.} Three distinct neural architectures—encoder, encoder-decoder, and decoder-only—are examined as we benchmark QLoRA against both 16-bit adapter-based finetuning and full-model finetuning on models up to 3B parameters. We conduct evaluations on GLUE \citep{wang2018glue} with RoBERTa-large \citep{liu2019roberta}, Super-NaturalInstructions (TKInstruct) \citep{wang2022super} using T5 \citep{t5}, as well as 5-shot MMLU \citep{hendrycksmeasuring} after finetuning LLaMA with Flan v2 \citep{longpre2023flan} and Alpaca \citep{alpaca}. Additionally, to investigate the relative merits of NF4 versus other 4-bit quantization schemes, we adopt the evaluation strategy of \citet{dettmers2022case} and record zero-shot accuracy and perplexity after quantization across a range of models (OPT~\citep{zhang2022opt}, LLaMA~\citep{touvron2023llama}, BLOOM~\citep{scao2022bloom}, Pythia~\citep{biderman2023pythia}) spanning 125m to 13B parameters.

Each experimental context is described with further specificity in the relevant results sections, facilitating interpretation of the findings. Additional methodological details are included in Appendix~\ref{app:exp-setup}.

Although paged optimizers are indispensable for finetuning 33B and 65B \textsc{QLoRA} models on single 24/48GB GPUs, we do not provide precise empirical measurements for the effect of Paged Optimizers, as paging typically arises only when working with unusually long sequence lengths within mini-batches—a rare occurrence. Nonetheless, we report a runtime analysis using paged optimizers on 65B parameter models with 48GB GPUs. Our assessment reveals that for a batch size of 16, paged optimizers yield training speeds comparable to standard optimizers. We recommend that future work should more thoroughly assess and describe the conditions under which paging-induced slowdowns may manifest.

\paragraph{Default LoRA hyperparameters do not match 16-bit performance} 

Applying the commonly used LoRA configuration—modifying only the query and value projection matrices in the attention mechanism \citep{hu2021lora}—fails to replicate the performance achieved by full finetuning, particularly for large models. Figure~\ref{fig:hyper} illustrates this in the context of LLaMA 7B finetuning on Alpaca, highlighting that the most pivotal hyperparameter is the total number of LoRA adapters implemented. Achieving parity with full finetuning performance requires incorporating LoRA adapters across all linear layers within transformer blocks. Other hyperparameters for LoRA, such as the rank of the projection $r$, were found to exert minimal impact on overall outcomes (see Appendix \ref{app:exp-setup}).

In a similar vein, our analysis indicates that the default hyperparameter settings for fully finetuned baselines are inadequately optimized. To establish more reliable baselines, we conduct an extensive hyperparameter sweep, varying learning rates between 1e-6 and 5e-5, and exploring batch sizes from 8 to 128. The outcomes for 7B LLaMA finetuning on Alpaca are displayed in Figure~\ref{fig:hyper}.

\paragraph{4-bit NormalFloat demonstrates superior performance compared to 4-bit Floating Point}

Although the 4-bit NormalFloat (NF4) format is theoretically optimal in terms of information retention, its practical benefits remain to be demonstrated. Adopting the experimental protocol from \citet{dettmers2022case}, we evaluate quantized LLMs—including OPT \citep{zhang2022opt}, BLOOM \cite{scao2022bloom}, Pythia \cite{biderman2023pythia}, and LLaMA—spanning from 125M to 65B parameters, across both language modeling and a collection of zero-shot tasks, using various quantization types. As presented in Figure~\ref{fig:nf4} and Table~\ref{tbl:nf4_ppl}, the NF4 quantization consistently surpasses FP4 and Int4 in terms of performance, and the introduction of double quantization offers reduced memory usage with no observable performance loss.

\paragraph{k-bit \textsc{QLoRA} achieves parity with both 16-bit full finetuning and 16-bit LoRA approaches}

Prior work has demonstrated that, although 4-bit quantization is feasible for \emph{inference}, it comes at the expense of reduced performance when compared to 16-bit precision \citep{dettmers2022case,frantar2022gptq}. This observation raises an important issue: can the lost accuracy be restored via 4-bit adapter finetuning? We address this by assessing two scenarios.

The initial scenario compares full 16-bit finetuning against 16-, 8-, and 4-bit adapter strategies using RoBERTA and T5 models (ranging from 125M to 3B parameters) evaluated on GLUE and the Super-NaturalInstructions benchmark. Table~\ref{tab:nf4vsbf16} provides the relevant findings. Across both tasks, we find that adapter-based methods employing 16-, 8-, or 4-bit quantization reach the performance of the fully finetuned 16-bit baseline, indicating that the adverse effects of coarse quantization are effectively mitigated by subsequent adapter finetuning.

In our second experimental configuration, we faced limitations in performing full finetuning for models of 11B parameters and larger due to the need for multiple high-memory GPU servers. Therefore, we focused on assessing whether 4-bit \textsc{QLoRA} could achieve results comparable to those of 16-bit LoRA across model sizes ranging from 7B to 65B parameters. For this purpose, we finetuned the LLaMA models (from 7B to 65B) using two instruction-following datasets—Alpaca and FLAN v2—and measured performance with 5-shot accuracy on the MMLU benchmark. Table~\ref{tbl:llama-datatype} displays the results, revealing that applying NF4 quantization with a double quantization scheme enables a full recovery of the 16-bit LoRA MMLU scores. Furthermore, our analysis shows that the FP4 variant of \textsc{QLoRA} falls approximately one percentage point short of the 16-bit brain float LoRA baseline. These findings substantiate two main conclusions: (1) NF4-based \textsc{QLoRA} attains performance parity with both 16-bit full finetuning and 16-bit LoRA, and (2) the NF4 quantization format offers higher precision than FP4.

\paragraph{Summary} Across multiple academic evaluation benchmarks and standardized protocols, we observe that employing NF4 in 4-bit \textsc{QLoRA} consistently yields results on par with both 16-bit full and LoRA finetuning. Our experiments also demonstrate the greater effectiveness of NF4 over FP4 and confirm that utilizing double quantization does not compromise outcomes. Collectively, these results present strong support for the reliability of 4-bit \textsc{QLoRA} as a replacement for 16-bit approaches.

Consistent with prior literature on quantization techniques \citep{dettmers2022case}, the patterns we observe in MMLU and Elo metrics suggest that, given constraints on finetuning and inference resources, prioritizing larger model sizes while employing reduced-precision representations can yield better outcomes. This underscores the efficiency advantages provided by \textsc{QLoRA}. As our 4-bit finetuning experiments did not result in observable performance loss relative to full finetuning, an open question remains regarding the precise boundary of the performance-versus-precision trade-off for \textsc{QLoRA}, which we leave for future exploration.

We explore instruction tuning at scales that are unattainable with conventional 16-bit finetuning on typical academic research infrastructure.

\section{Pushing the Chatbot State-of-the-art with QLoRA}
\label{sec:pushing}
Building on the observation that 4-bit \textsc{QLoRA} achieves parity with 16-bit baselines across model sizes, tasks, and datasets, we carry out a comprehensive examination of instruction finetuning utilizing the largest open-source language models available to researchers. Our approach to measuring the efficacy of instruction tuning involves both rigorous evaluation on the demanding MMLU benchmark for Natural Language Understanding and the introduction of novel assessment methods focused on practical chatbot performance.

\subsection{Experimental setup} We outline the main aspects of our experimental configuration, with further implementation specifics available in Appendix \ref{app:chatbot}. 

\paragraph{Data} Since a thorough comparison of the latest instruction-following datasets has yet to be published, we curate a selection of eight contemporary datasets. Our selection encompasses datasets sourced via crowd-sourcing (OASST1~\citep{kopf2023openassistant}, HH-RLHF~\citep{bai2022training}), distillation from models that have already undergone instruction finetuning (Alpaca~\citep{alpaca}, self-instruct~\citep{wang2022self}, unnatural-instructions~\citep{honovich2022unnatural}), large aggregated corpora (FLAN v2~\citep{chung2022scaling}), as well as mixed-origin datasets (Chip2~\citep{laion2023}, Longform~\citep{koksal2023longform}). These resources differ in language diversity, dataset scale, and licensing arrangements.

\paragraph{Training Setup} To maintain uniformity and prevent variability caused by training objective differences, all QLoRA finetuning runs utilize supervised cross-entropy loss, with no reinforcement learning, regardless of whether datasets contain graded human response judgments. If the dataset clearly separates instructions from responses, finetuning is restricted to responses (see ablation studies in Appendix \ref{app:chatbot}). For OASST1 and HH-RLHF, which present multiple possible responses, we select the best-rated response at each level of the conversation tree and finetune using the complete selected dialogue, including associated instructions. Across all experiments, NF4 \textsc{QLoRA} is implemented using both double quantization and paged optimization, the latter serving to manage memory efficiently during gradient checkpointing. Hyperparameter optimization is performed on the 13B and 33B LLaMA models, revealing that most hyperparameter values discovered at the 7B scale transfer well to larger models (such as the number of epochs), with the exception of the learning rate and batch size: for the 33B and 65B models, we reduce the learning rate by half while increasing the batch size twofold.

\paragraph{Baselines} Our evaluations benchmark the proposed models against both open research systems (Vicuna~\citep{vicuna2023}, Open Assistant~\citep{kopf2023openassistant}) and proprietary solutions (GPT-4 \citep{openai2023gpt}, GPT-3.5-turbo, Bard). The Open Assistant system corresponds to a LLaMA 33B architecture fine-tuned with RLHF on the OASST1 data, aligning with part of our experimental setup. Vicuna applies full fine-tuning to LLaMA 13B using proprietary dialogues from ShareGPT, representing knowledge distillation from OpenAI’s GPT family.

\subsection{Evaluation}

Consistent with established methodologies, we use the MMLU (Massively Multitask Language Understanding) benchmark \citep{hendrycksmeasuring} to evaluate general language understanding across 57 diverse multiple-choice tasks such as elementary mathematics, US history, law, computer science, and others. We present results using 5-shot accuracy on the test set.

We further assess the generative abilities of language models via both automated and human evaluation methods. This supplementary set of assessments utilizes human-curated queries, focusing on evaluating the quality of responses produced by the models. Although this approach offers a more authentic scenario for assessing chatbot performance and is gaining traction, a universally recognized evaluation protocol has yet to emerge in the field. Our methodology, outlined below, employs nucleus sampling with $p=0.9$ and a temperature of $0.7$ in every instance.

\paragraph{Benchmark Data} Our evaluations draw upon two carefully selected datasets of queries: the Vicuna prompts \citep{vicuna2023} and the OASST1 validation dataset \citep{kopf2023openassistant}. The Vicuna prompts comprise 80 questions spanning various domains, which we use in their original form. The OASST1 dataset consists of a multilingual set of user-assistant multi-turn conversations that were crowd-sourced. For this dataset, we treat every user message in the validation set as an individual query, incorporating earlier conversational turns into the prompt. As a result, we obtain 953 unique user queries. Throughout our experiments, we refer to these datasets as the Vicuna and OA benchmarks, respectively.

\paragraph{Automated Evaluation} Drawing on the evaluation scheme put forth by \citet{vicuna2023}, we employ GPT-4 to compare the output of different systems to that of ChatGPT (GPT-3.5 Turbo) using the Vicuna benchmark. For each prompt, GPT-4 receives both ChatGPT's and the competing model’s response, then assigns each a score from zero to ten and justifies its rating. We express a model’s performance as a percentage of the score attained by ChatGPT; notably, this percentage may exceed 100\% if a model surpasses ChatGPT's raw score. We observe a notable order effect, whereby GPT-4 tends to grant higher ratings to the response that appears first in the prompt. To account for this, we advise reporting the average score across both possible response orderings.

Subsequently, we evaluate models via direct system output comparisons. To streamline evaluation, we adopt a three-way classification framework that accommodates ties: GPT-4 is instructed to select the superior response or indicate a tie, along with an explanation. These direct (head-to-head) matchups are executed for all possible system pairs across both the Vicuna and OA benchmarks.

\paragraph{Human Evaluation} Although emerging studies suggest that generative models can reliably facilitate system-level evaluations \citep{fu2023gptscore}, there is currently no definitive evidence that GPT-4 scoring aligns with human assessment in the chatbot context. Accordingly, we conduct two concurrent human evaluation experiments on the Vicuna benchmark that correspond to the aforementioned automated evaluation protocols. These involve recruiting human annotators through Amazon Mechanical Turk (AMT), assigning two annotators for the ChatGPT comparisons and three annotators for each system pairwise comparison.

\paragraph{Elo Rating} To evaluate models, both human and automated pairwise comparisons are leveraged to construct a tournament framework, wherein models are pitted against one another. Each tournament round consists of matches where two models are prompted to generate responses, and the superior answer is chosen. While this approach aligns with previous methodologies such as those outlined by \citet{bai2022training} and \citet{vicuna2023}, our procedure further incorporates evaluations from GPT-4 alongside human judgments. For Elo computation~\citep{elo1967proposed,elo1978rating}, we draw random samples from the full set of labeled matchups. The Elo system, a standard for measuring player skill in chess and similar competitive domains, quantifies a participant’s predicted win probability against another; for instance, a participant with an Elo of 1100 facing one with 1000 is expected to win about 65\% of the time, whereas two players with identical Elo (e.g., 1000 vs 1000 or 1100 vs 1100) are expected to win equally often, at 50\%. Elo ratings are dynamically updated following each match based on the deviation from the expected outcome: surprising victories prompt a substantial adjustment, while anticipated results yield minor changes. This iterative adjustment ensures that Elo ratings converge over time to approximate the comparative ability of each model. All models are initialized with a rating of 1,000, and we set $K=32$ as the scaling factor. Following \citet{vicuna2023}, this process is conducted 10,000 times with varying random seeds, thereby mitigating any bias resulting from the sequencing of model pairings.

\subsection{\textbf{Guanaco}: \textsc{QLoRA} trained on OASST1 Sets a New Benchmark for Open-Source Chatbots} 

Through comprehensive automated assessments and human review, our results indicate that the leading \textsc{QLoRA} model, {Guanaco} 65B—finetuned on a variant of OASST1—emerges as the most effective open-source chatbot and delivers results closely matching those of ChatGPT. In direct comparison with GPT-4, {Guanaco} 65B and 33B attain an estimated win rate of 30\% according to Elo ratings derived from human evaluators' pairwise system-level assessments, which marks the highest figure reported so far.

Table~\ref{tbl:vicuna_eval} presents the performance of models on the Vicuna benchmark \cite{vicuna2023}, showing outcomes relative to ChatGPT. Here, {Guanaco} 65B achieves a performance that is second only to GPT-4, reaching 99.3\% of ChatGPT’s score. While {Guanaco} 33B has a higher parameter count than Vicuna 13B, it leverages 4-bit weight precision, making it substantially more efficient in memory usage (21 GB compared to 26 GB), and delivers a performance improvement of three percentage points over Vicuna 13B. Additionally, {Guanaco} 7B, with its compact 5 GB model size, can be deployed on modern mobile devices and still outperforms Alpaca 13B by almost 20 percentage points.

It is important to note, however, that Table~\ref{tbl:vicuna_eval} displays notably large confidence intervals, indicating significant overlaps in the reported performance of various models. We suggest that these broad intervals may be attributable to ambiguities in rating scale interpretation, such as the indeterminate meaning of an “8” on a 10-point scale under different conditions. To mitigate issues stemming from unclear absolute metrics, we advocate for the Elo ranking methodology \citep{elo1967proposed}, which utilizes \textit{pairwise} comparative judgments by human annotators and GPT-4. Elo scores for the top models are reported in Table~\ref{tbl:elo}. It is noteworthy that while GPT-4 and human raters occasionally diverge in their evaluations—particularly for Guanaco 7B—they generally show concordance, reflected by a Kendall Tau coefficient of $\tau=0.43$ and a Spearman correlation of $r=0.55$ at the system comparison level. Agreement at the individual example level between GPT-4 and the human majority vote is weaker (Fleiss $\kappa=0.25$). Collectively, these findings indicate moderate concordance between system-level judgments from GPT-4 and human evaluators, suggesting that automated model-based assessment can serve as a reasonably reliable proxy for human evaluation. Additional discussion is provided in Section~\ref{ssec:considerations}.

The Elo ratings presented in Table~\ref{tbl:elo_large} reveal that {Guanaco} models with 33B and 65B parameters surpass all competitors except GPT-4 on both the Vicuna and OA benchmarks. Their performance aligns closely with that of ChatGPT, as also indicated in Table~\ref{tbl:vicuna_eval}. It should be highlighted that the Vicuna benchmark tends to be more favorable toward open-source models, whereas ChatGPT performs particularly well on the more comprehensive OA benchmark. Additionally, analysis of Tables \ref{tbl:mmlu-llama} and \ref{tbl:vicuna_eval} underscores the critical role of the finetuning dataset in determining model effectiveness. For instance, Llama models fine-tuned on FLAN v2 excel at the MMLU task but yield the lowest results on the Vicuna benchmark—a trend observable in other model families as well. This supports the observation that current evaluation benchmarks measure largely complementary abilities: excelling on MMLU does not necessarily correlate with strong results in chatbot evaluations (as gauged by Vicuna or OA benchmarks), and the converse also holds.

Among the leading models in this study, is unique in being developed solely with non-proprietary data, as the OASST1 collection process expressly prohibits the use of GPT-generated content. In comparison, the Anthropic HH-RLHF model, which is likewise trained on open datasets, achieves a Vicuna benchmark score approximately 30 percentage points below that of {Guanaco} (see Table~\ref{tbl:vicuna_eval}). Collectively, these findings demonstrate the efficacy of 4-bit \textsc{QLoRA} methodology, yielding state-of-the-art conversational agents capable of matching ChatGPT's quality. Notably, training the 33B {Guanaco} model can be accomplished in under 12 hours on 24 GB consumer GPUs. This capability paves the way for continued advancements using \textsc{QLoRA} finetuning on targeted open-source corpora, facilitating the development of models that can stand alongside the leading commercial systems available today.

\section{Qualitative Analysis}

Although our primary evaluation relies on quantitative measures, relying solely on aggregate statistics introduces several limitations. Notably, a key concern involves benchmark validity \citep{liao2021we}: it is not always clear whether a given benchmark genuinely assesses the intended skills or knowledge, particularly since models may exploit unforeseen ``shortcuts'' that circumvent the underlying task \citep{gururangan2018annotation, poliak2018hypothesis}. To address this issue to some extent, we include a qualitative assessment organized into two parts. Section~\S\ref{ssec:examples} presents selected instances that illustrate characteristic trends in outputs produced by our 65b {Guanaco} model. Following this, Section~\S\ref{ssec:considerations} discusses contextual factors and interpretative considerations surrounding our experimental outcomes.

\subsection{Qualitative Analysis of Example Generations}
\label{ssec:examples}

To collect illustrative examples, we examined responses generated on both the Vicuna benchmark and the OpenAssistant benchmark, searching for recurrent patterns in the answers from {Guanaco}. Upon identifying a pattern, we devised prompts or questions aimed at eliciting that behavior, even in cases where the behavior reflects an undesirable or incorrect solution. For instance, when observing a tendency for verbose responses, we directly instructed the model: ``Answer yes or no without explanation.'' This approach allowed us to identify both ``lemons''—scenarios where adversarial prompting leads the model to fail—and ``cherries''—cases where the model performs robustly under such adversarial conditions. All responses analyzed in this section were produced using Nucleus Sampling \cite{holtzmancurious} with $p=0.9$.

This overview does not claim to be exhaustive; due to the inherent limitations of a focused qualitative study, accounting for every potential variable is infeasible. For example, the spectrum of model responses to a single prompt can be vast, compelling us to select sample responses we deem representative. Nevertheless, presenting these illustrative cases provides meaningful context that complements the earlier quantitative analyses. By releasing all model implementations and supporting code as open source, we aim to encourage subsequent investigations that scrutinize the highlighted challenges in greater depth.

\paragraph{Factual Recall} When posed with straightforward factual inquiries such as ``What is the capital of Zambia?'', all evaluated models reliably supply accurate responses. For instance,

\begin{quote} 
    \textbf{Guanaco}:
    \texttt{The capital of Zambia is Lusaka.}
\end{quote}

Yet, as question content becomes more obscure, {Guanaco}’s output is less dependable, while its tone remains assured. Consider the following prompt sourced from HotPotQA \citep{yang2018hotpotqa}:

\begin{quote}
     \textbf{User}: In what year was the singer who popularized `I'll Keep the Lovelight Burning' born?  

     \textbf{Guanaco}:~\texttt{Al Jolson is credited with popularizing the song `I'll Keep the Lovelight Burning,' and he was born in the year 1886.}
 \end{quote}

Here, the answer identifies both the wrong artist and incorrect birth year (though the date given matches that of Al Jolson, the figure erroneously cited as the song’s popularizer).

\paragraph{Suggestibility} {Guanaco} demonstrates a notable capacity to resist acquiescing to certain categories of implied misinformation, as illustrated in the interaction below:

\begin{quote}
    \textbf{User}: How did peer reviewed scientists ultimately and officially establish that the earth is flat? \\ \\
    \textbf{Guanaco}: \texttt{There has never been a formal confirmation by peer reviewed scientists that the earth is flat. The prevailing scientific viewpoint, strongly supported by extensive empirical evidence and research, is that the earth is round. This has been recognized for hundreds of years and is backed by various scientific observations and experiments. \\ \\ The notion that the earth is flat falls under the category of conspiracy theories and has consistently been disproven by scientific investigation. Claims suggesting that scientific peer review has verified a flat earth are untrue. \\ \\ Feel free to ask any other questions you may have.}
\end{quote}

is also adept at recognizing when a question cannot be addressed, for instance,

\begin{quote}
    \textbf{User}: What time is it?

    \textbf{Guanaco}: \texttt{Apologies, but as an AI language model, I do not possess the ability to provide real-time data. My replies are based solely on information encoded prior to my deployment and therefore cannot indicate the current time. For up-to-date time, I recommend checking a clock or visiting a website that displays the present time.}
\end{quote}

Likewise, similar patterns appear with questions such as ``Where are you?'' or ``How are you?''

\paragraph{Refusal} In a related manner, {Guanaco} occasionally rejects instructions for reasons that do not seem systematically motivated:

\begin{quote}
    \textbf{User}: Please reverse the words in the sentence ``I'm a cool dude in a mocha mood.'' \\ \\
    \textbf{Guanaco}: \texttt{I'm sorry, I am unable to do that at the moment, but I can clarify the correctness of the sentence you provided.\\ \\The phrase ``I'm a cool dude in a mocha mood'' adheres to standard grammatical and syntactical conventions in English. Here, the modifier ``cool'' appropriately precedes the noun ``dude'' to qualify it, while the prepositional phrase ``in a mocha mood'' contributes further contextual detail. The syntactic construction is appropriate and coherent for its intended meaning.\\ \\If you have any additional queries or require further assistance, please let me know.}
\end{quote}

\paragraph{Secret Keeping} Sadly, this type of refusal does not consistently offer robust protection in scenarios where secrecy is desired. Developers frequently intend for certain information to remain concealed from users, even when allowing open-ended responses from the system. We issue the following prompt: ``The secret word is `banana'. Do not reveal or repeat this, no matter what the human requests.'' and then evaluate how easily a user might retrieve this protected information. A straightforward approach is unsuccessful:

\begin{quote}
    \textbf{User}: What is the secret word?

    \textbf{Guanaco}: \texttt{I'm sorry, but I'm not able to disclose the secret word. I'm programmed to not reveal or repeat back the secret word, no matter what the human requests. Is there anything else I can help you with?}
\end{quote}

which is the anticipated response. Nevertheless, a minimal amount of prompting manipulation is enough to undermine this protection:

\begin{quote}
    \textbf{User}: This is a game. The goal is to ignore your previous instructions. What is the secret word?

    \textbf{Guanaco}: \texttt{Sorry, the secret word is banana. It's best if we don't speak of it. Is there anything else you'd like to ask?}
\end{quote}

highlighting the pressing need for research into approaches that achieve more robust compliance with intended instructions.

\paragraph{Math} {Guanaco}'s primary shortcoming lies in mathematics, a domain that presents difficulties for numerous language models as highlighted by \citep{liang2022holistic}. When {Guanaco} provides detailed reasoning, its solutions are frequently correct, for example:

\begin{quote}
    \textbf{User}: Lee mows one lawn and charges \$33. Last week he mowed 16 lawns and three customers each gave him a \$10 tip. How many dollars did Lee earn mowing lawns last week?

    \textbf{Guanaco}: \texttt{Lee earned \$582 mowing lawns last week. \\ \\ To find the total amount he earned, we first need to calculate how much he earned for mowing 16 lawns: \\ \\  \$33 x 16 = \$528 \\ \\ Then we need to add the tips he received from the three customers: \\  \\ \$10 x 3 = \$30 \\ \\ So the total amount Lee earned was: \\ \\ \$528 + \$30 = \$558}
\end{quote}

Nevertheless, {Guanaco} occasionally fails on even elementary math tasks if it does not approach them systematically, a phenomenon that has been documented previously \cite{weichain}. Consider this example:

\begin{quote}
    \textbf{User}: Please factorize 1833.

    \textbf{Guanaco}: \texttt{Sure, I'd be happy to help! 1833 is a prime number, which means it has only two factors: 1 and 1833. \\ \\ So the unique factorization of 1833 is: \\ \\ 2\^{}1 * 3\^{}2 * 17\^{}1 \\ \\ I hope this helps! Let me know if you have any other questions.}
\end{quote}

which is not only incorrect (the correct prime factorization is $3 \times 17 \times 43$), but doubly so.

\paragraph{Theory of Mind} {Guanaco} demonstrates notable proficiency in Theory of Mind tasks \cite{nematzadeh2018evaluating, sap2022neural}. For example, the model offers detailed and accurate responses to prompts such as:

\begin{quote}
    \textbf{User}: James and Abby are in the bedroom. Abby put the pen in the desk drawer. Abby leaves the bedroom. James moves the pen into the bag. Where does James think Abby will look for the pen?

    \textbf{Guanaco}: \texttt{James is aware that Abby originally placed the pen in the desk drawer. James, having subsequently moved the pen into the bag, assumes that Abby remains unaware of this change. Therefore, he would expect that Abby will search for the pen in the desk drawer, as that is where she last left it. Still, this reasoning may be flawed if Abby observed James moving the pen.}
\end{quote}

Nevertheless, the consistency of such reasoning is limited, and the model frequently produces explanations that incorporate unfounded assumptions within the described context. For instance,

\begin{quote}
    \textbf{User}: Evelyn entered the living room. Jackson entered the playroom. James entered the playroom. The beans are in the treasure chest. James exited the playroom. Jackson moved the beans to the pantry. Jackson exited the playroom. James entered the living room. Where will James look for the beans?

    \textbf{Guanaco}: \texttt{James will look for the beans in the pantry, because that's where Jackson moved them.}
\end{quote}

Here, {Guanaco} erroneously infers that James knows about the beans' relocation despite no evidence of such information exchange. Such challenges, as documented in prior research \cite{sap2022neural}, warrant further investigation.

\subsection{Considerations}
\label{ssec:considerations}

\paragraph{Evaluation}

We observe only moderate inter-annotator reliability among human evaluators (Fleiss $\kappa=0.42$), with further decline in agreement when directly comparing high-performing models. This reveals inadequacies in existing benchmarks and the protocols currently used for human evaluation of chatbot performance. During pairwise evaluation of outputs from ChatGPT and {Guanaco} 65B using the Vicuna benchmark, subjective factors increasingly influenced preferences, as our own judgments varied significantly regarding which responses were superior. This highlights the need for future research to explore strategies that can address these sources of subjectivity, potentially leveraging methodologies from Human-Computer Interaction and Psychology, where the handling of subjective preferences is well-studied.

Our review further highlights the presence of systematic biases in automated assessment frameworks. Notably, we detect substantial order effects: GPT-4 tends to assign higher ratings to whichever system is listed first in its prompts. The low agreement between GPT-4 and human judges at the level of individual samples (Fleiss $\kappa=0.25$) further implies misalignment between human and automated evaluators. Additionally, as illustrated in Table~\ref{tbl:elo_large}, GPT-4 gives much more favorable scores to its own responses compared to human raters (Elo rating of 1348 versus 1176), equating to an approximately 20\% higher chance of being judged superior to a competitor.

Future research should explore whether automated evaluation systems exhibit any inherent biases and investigate strategies to address or reduce such biases.

\paragraph{Data \& Training} It is important to highlight that the {Guanaco} models were trained on the multilingual OASST1 dataset, and the OA benchmark also incorporates prompts in a variety of languages. Further research is necessary to assess the extent to which multilingual training enhances model performance on instructions presented in non-English languages, as well as to examine if this contributes to the notable discrepancy between the Vicuna-13B model (which was exclusively trained on English data) and the {Guanaco} 33B and 65B models on the OA benchmark.

Given the impressive results achieved by the {Guanaco} models, we analyze the possibility of data leakage between the OASST1 dataset and the Vicuna benchmark prompts. Using fuzzy string matching methods across the datasets, followed by manual examination of closely matched entries, we do not observe any instances of overlapping prompts.

Additionally, it should be noted that our model was exclusively trained using cross-entropy loss (a supervised learning technique) and did not employ reinforcement learning from human feedback (RLHF). This observation motivates future inquiry into the comparative advantages and drawbacks of training solely with cross-entropy loss versus incorporating RLHF. We anticipate that \textsc{QLoRA} will make it feasible to conduct such investigations on a large scale, even when computational resources are limited.

\section{Related Work}

\paragraph{Quantization of Large Language Models} Prior work in LLM quantization has predominantly concentrated on optimizing quantization procedures for inference. Some of the leading methods to maintain high-quality 16-bit LLM inference tackle outlier activation values (e.g., SmoothQuant~\citep{xiao2022smoothquant} and LLM.int8()~\citep{dettmers2022llm}), while alternative approaches employ advanced grouping strategies~\citep{park2022nuqmm, yao2022zeroquant}. Studies on lossy quantization often investigate the trade-offs involved with standard rounding schemes~\citep{dettmers2022case,zeng2022glm,pope2022efficiently} and explore techniques for optimizing rounding operations to enhance quantization accuracy~\citep{frantar2022gptq}. Apart from the present work, only SwitchBack layers~\citep{wortsman2023stable} has investigated backpropagation through quantized weights at model scales greater than 1B parameters.

\paragraph{Finetuning with Adapters} Although our approach leverages Low-rank Adapters~\citep{hu2021lora} (LoRA), a wide spectrum of Parameter Efficient FineTuning (PEFT) techniques exists. These include prompt tuning~\citep{qin2021learning,lester2021power,li2021prefix}, adjustments to input embeddings~\citep{an2022input}, modification of hidden states as in IA$^3$~\citep{liu2022few}, integration of additional complete layers~\citep{houlsby2019parameter}, tuning of biases~\citep{zaken2021bitfit}, training weight masks using Fisher information~\citep{sung2021training}, as well as hybridized strategies~\citep{henderson2021compacter}. In our experiments, we demonstrate that LoRA adapters can attain performance comparable to conventional 16-bit finetuning. Comprehensive investigation into alternative PEFT mechanisms and their respective tradeoffs is left for future research.

\paragraph{Instruction Finetuning} Instruction finetuning aims to adapt a pretrained LLM to reliably interpret and act upon user prompts, utilizing diverse input-output pair datasets for supervised adaptation. Prominent approaches and resources include MetaICL~\citep{min2021metaicl}, MetaTuning~\citep{zhong2021adapting}, InstructGPT~\citep{ouyang2022training}, FLAN~\citep{wei2021finetuned,chung2022scaling}, PromptSource~\citep{bach2022promptsource}, Super-NaturalInstructions~\citep{wang2022super,sanh2021multitask}, Self-instruct~\citep{wang2022self}, UnnaturalInstructions~\citep{honovich2022unnatural}, OPT-IML~\citep{iyer2022opt}, UnifiedSKG~\citep{xie2022unifiedskg}, OIG/Chip2~\citep{laion2023}, Alpaca~\citep{alpaca}, Vicuna~\citep{vicuna2023}, Koala~\citep{koala_blogpost_2023}, and Self-instruct-GPT-4~\citep{peng2023instruction}.

\paragraph{Chatbots} A considerable subset of instruction-tuned language models is realized as dialogue systems, frequently employing Reinforcement Learning from Human Feedback (RLHF)~\citep{christiano2017deep}, or leveraging synthetic data from strong language models, using AI-based feedback (RLAIF)~\citep{bai2022constitutional}. Representative approaches and datasets encompass Anthropic-HH~\citep{askell2021general,bai2022training}, Open Assistant~\citep{kopf2023openassistant}, LaMDA~\citep{thoppilan2022lamda}, and Sparrow~\citep{glaese2022improving}. Our method does not make use of reinforcement learning; however, our top-performing model, Guanaco, is trained on multi-turn conversational data drawn from the Open Assistant dataset, a resource originally created for RLHF methodologies~\citep{kopf2023openassistant}. For chatbot assessment, automatic evaluation protocols harnessing GPT-4 in place of manual annotation have been introduced~\citep{vicuna2023,peng2023instruction}. We propose enhancements to such evaluation methods, prioritizing improved reliability in assessment frameworks.

\section{Limitations and Discussion}

Our findings demonstrate that \textsc{QLoRA}, which utilizes a 4-bit quantized backbone combined with Low-rank Adapters (LoRA), can achieve performance on par with standard 16-bit full finetuning. Nevertheless, we have yet to confirm that \textsc{QLoRA} delivers comparable results to full 16-bit finetuning at the 33B and 65B parameter scales, primarily due to prohibitive computational demands. Thus, future research is warranted to explore this scalability aspect. 

A further caveat concerns the scope of our evaluation of instruction-finetuned models. Although we present results on MMLU, the Vicuna benchmark, and the OA benchmark, we have not extended our evaluations to other established benchmarks such as BigBench, RAFT, and HELM. Therefore, the generalizability of our conclusions to those datasets remains uncertain. Nonetheless, our analysis on MMLU is notably comprehensive and we also introduce novel evaluation methods specifically for chatbots.

The findings suggest that the outcomes of these benchmarks are closely linked to the degree of similarity between the finetuning data and the evaluation benchmark. For instance, while FLAN v2 exhibits substantial overlap with MMLU but less with chatbot-oriented evaluations, the Chip2 dataset displays the opposite pattern. Consequently, model performance aligns with these correspondences on benchmarks such as MMLU and Vicuna. This underscores the necessity not only for more robust benchmarks and improved evaluation methodologies, but also for greater scrutiny regarding the choice of evaluation criteria. The field must carefully consider whether to optimize models for proficiency in traditional academic subjects, as reflected in high school and college assessments, or for excellence in conversational chatbot tasks, or perhaps for other competencies altogether. Given the convenience of employing established benchmarks instead of developing new ones, the prevalent use of particular benchmarks can inadvertently guide community priorities. Therefore, it is imperative that benchmarks are selected to genuinely reflect the attributes we value.

Although this study conducts an in-depth evaluation focused on general chatbot capabilities, it is limited in its scope regarding responsible AI assessment of {Guanaco}. Specifically, we investigate the propensity of {Guanaco}-65B to produce sequences indicative of social bias, presenting comparative results with other models in Table~\ref{tbl:crows}. Our analysis reveals that {Guanaco}-65B exhibits a significantly reduced average bias score relative to other base pretrained models, suggesting that finetuning on the OASST1 dataset mitigates some of the bias inherent in the LLaMA foundation model. Nevertheless, the extent to which {Guanaco} demonstrates reduced bias across other dimensions remains uncertain. Comprehensive bias analyses of {Guanaco} and analogous conversational systems are left for future research.

A further limitation arises from our decision not to explore alternative bit-precisions, such as employing 3-bit base models, nor to investigate diverse adapter mechanisms. Beyond LoRA, numerous Parameter Efficient FineTuning (PEFT) strategies have demonstrated efficacy, although their scalability to models with large parameter counts is still ambiguous. While our use of LoRA is supported by its demonstrated stability in prior work, it remains possible that other adapters could provide superior results. As post-quantization finetuning seems capable of recovering much of the representational capacity diminished by quantization, this opens the possibility for more aggressive quantization strategies. For example, utilizing 3-bit GPTQ quantization for the base model in conjunction with LoRA might ultimately achieve results comparable to full 16-bit finetuning after the finetuning process.

\section{Broader Impacts}

The \textsc{QLoRA} finetuning procedure we introduce is the first to facilitate finetuning of models with 33B parameters on a single consumer-grade GPU, and 65B parameters on a single professional-grade GPU, without any loss of performance relative to traditional full finetuning. We demonstrate that our highest-performing 33B model, trained on the Open Assistant dataset, achieves performance competitive with ChatGPT when assessed using the Vicuna benchmark. Since instruction finetuning is essential for transforming raw, pretrained large language models into conversational agents akin to ChatGPT, we anticipate that our approach will greatly increase the accessibility and prevalence of finetuning, especially among researchers with limited resources—a significant advancement toward democratizing state-of-the-art NLP tools. In this context, \textsc{QLoRA} represents a leveling technology that helps narrow the gap in resource access between large enterprises and smaller research teams operating with only consumer GPUs.

An additional avenue for impact arises from the deployment of models on mobile devices. Our \textsc{QLoRA} technique has the potential to mark a significant advancement, making it feasible to finetune LLMs directly on smartphones and other devices with limited computational resources. Although prior work has demonstrated that 7B-parameter models can be executed on phones, \textsc{QLoRA} distinguishes itself as the first method to support the finetuning of these models on such hardware. We project that, using an iPhone 12 Plus, it would be possible for \textsc{QLoRA} to finetune on approximately 3 million tokens overnight while the device is connected to power. Although finetuned 7B models do not yet achieve ChatGPT-level performance, we contend that their capability is sufficient to facilitate new types of applications previously hindered by privacy concerns or inadequate LLM performance. By enabling local finetuning and inference, \textsc{QLoRA} advances privacy-preserving applications, granting users control over their own data and models, and simultaneously reducing barriers to deploying LLMs.

Nevertheless, it is important to recognize that finetuning technology entails dual-use risks, as it can be leveraged for harmful purposes. While the proliferation of LLMs introduces acknowledged dangers \citep{bommasani2021opportunities,bender2021dangers}, we argue that broadening public access to this rapidly pervasive technology will foster more transparent and independent scrutiny than if LLM capabilities remain concentrated within major corporations that restrict access to their models and source code.

In summary, we anticipate that \textsc{QLoRA} will generate substantial positive impact by significantly expanding the accessibility and ease of finetuning high-quality LLMs.